% Options for packages loaded elsewhere
\PassOptionsToPackage{unicode}{hyperref}
\PassOptionsToPackage{hyphens}{url}
\documentclass[
]{article}
\usepackage{xcolor}
\usepackage{amsmath,amssymb}
\setcounter{secnumdepth}{-\maxdimen} % remove section numbering
\usepackage{iftex}
\ifPDFTeX
  \usepackage[T1]{fontenc}
  \usepackage[utf8]{inputenc}
  \usepackage{textcomp} % provide euro and other symbols
\else % if luatex or xetex
  \usepackage{unicode-math} % this also loads fontspec
  \defaultfontfeatures{Scale=MatchLowercase}
  \defaultfontfeatures[\rmfamily]{Ligatures=TeX,Scale=1}
\fi
\usepackage{lmodern}
\ifPDFTeX\else
  % xetex/luatex font selection
\fi
% Use upquote if available, for straight quotes in verbatim environments
\IfFileExists{upquote.sty}{\usepackage{upquote}}{}
\IfFileExists{microtype.sty}{% use microtype if available
  \usepackage[]{microtype}
  \UseMicrotypeSet[protrusion]{basicmath} % disable protrusion for tt fonts
}{}
\makeatletter
\@ifundefined{KOMAClassName}{% if non-KOMA class
  \IfFileExists{parskip.sty}{%
    \usepackage{parskip}
  }{% else
    \setlength{\parindent}{0pt}
    \setlength{\parskip}{6pt plus 2pt minus 1pt}}
}{% if KOMA class
  \KOMAoptions{parskip=half}}
\makeatother
\usepackage{longtable,booktabs,array}
\newcounter{none} % for unnumbered tables
\usepackage{calc} % for calculating minipage widths
% Correct order of tables after \paragraph or \subparagraph
\usepackage{etoolbox}
\makeatletter
\patchcmd\longtable{\par}{\if@noskipsec\mbox{}\fi\par}{}{}
\makeatother
% Allow footnotes in longtable head/foot
\IfFileExists{footnotehyper.sty}{\usepackage{footnotehyper}}{\usepackage{footnote}}
\makesavenoteenv{longtable}
\setlength{\emergencystretch}{3em} % prevent overfull lines
\providecommand{\tightlist}{%
  \setlength{\itemsep}{0pt}\setlength{\parskip}{0pt}}
\usepackage[]{natbib}
\bibliographystyle{unsrtnat}
\usepackage{bookmark}
\IfFileExists{xurl.sty}{\usepackage{xurl}}{} % add URL line breaks if available
\urlstyle{same}
\hypersetup{
  pdftitle={dreamOfkiki: A Substrate-Agnostic Formal Framework for Dream-Based Knowledge Consolidation in Artificial Cognitive Systems},
  pdfauthor={dreamOfkiki project contributors},
  hidelinks,
  pdfcreator={LaTeX via pandoc}}

\title{dreamOfkiki: A Substrate-Agnostic Formal Framework for
Dream-Based Knowledge Consolidation in Artificial Cognitive Systems}
\author{dreamOfkiki project contributors}
\date{2026}

\begin{document}
\maketitle

\section{Paper 1 --- Full Draft
Assembly}\label{paper-1-full-draft-assembly}

⚠️ \textbf{Status} : draft assembly. Section .md files were the original
source of truth ; this file now inlines their content to become the
assembled source for pandoc rendering.

⚠️ \textbf{Synthetic data caveats} apply to §7 Results (numbers from
mega-v2 synthetic placeholder). Real ablation lands cycle 1 closeout
(S20+) or cycle 2.

\begin{center}\rule{0.5\linewidth}{0.5pt}\end{center}

\subsection{1. Abstract}\label{abstract}

Catastrophic forgetting remains a central obstacle for artificial
cognitive systems learning sequentially across tasks. Sleep-inspired
consolidation has been proposed as a remedy, with prior work exploring
replay (Walker, van de Ven), synaptic homeostasis (Tononi), creative
recombination (Hobson), and predictive coding (Friston) --- but no
unified formal framework integrates these four pillars into composable,
substrate-agnostic operations.

We introduce \textbf{dreamOfkiki}, a formal framework with executable
axioms (DR-0 accountability, DR-1 episodic conservation, DR-2
compositionality on a free semigroup of dream operations, DR-3
substrate-agnosticism via a Conformance Criterion, DR-4 profile chain
inclusion). The framework defines 8 typed primitives (α, β, γ, δ inputs
; 4 output channels), 4 canonical operations (replay, downscale,
restructure, recombine), and a 5-tuple Dream Episode ontology. The
framework admits multiple conformant substrates ; exemplar
implementations validate the design and are reported separately (see
Paper 2).

Pre-registered hypotheses (OSF DOI : pending) are evaluated via Welch's
t-test, TOST equivalence, Jonckheere-Terpstra trend, and one-sample
t-test against compute budget under Bonferroni correction.

\textbf{Pipeline Validation (synthetic placeholder, G2 pilot).} The
end-to-end measurement and statistical pipeline is exercised with mock
predictors at scripted accuracy levels ; numbers are reported in §7
alongside their registered run\_id and JSON dump under
\texttt{docs/milestones/}. Real mega-v2 inference and any fMRI
representational similarity analysis follow in cycle 2 (Paper 2). All
code, specifications, and pre-registration are open under MIT/CC-BY-4.0.

\begin{center}\rule{0.5\linewidth}{0.5pt}\end{center}

\subsection{2. Introduction}\label{introduction}

\subsubsection{2.1 Catastrophic forgetting and the consolidation
gap}\label{catastrophic-forgetting-and-the-consolidation-gap}

Modern artificial cognitive systems excel at single-task learning but
degrade rapidly when trained sequentially across tasks --- a phenomenon
known as \textbf{catastrophic forgetting}
\citep{mccloskey1989catastrophic, french1999catastrophic}. Despite two
decades of mitigation strategies (elastic weight consolidation
\citep{kirkpatrick2017overcoming}, generative replay
\citep{shin2017continual}, rehearsal-based memory
\citep{rebuffi2017icarl}), the field still lacks a \emph{unified
theoretical account} of why these mechanisms work and when they should
compose.

Biological cognition solves this problem during \textbf{sleep}.
Hippocampal replay during NREM, synaptic downscaling, predictive
restructuring of cortical representations, and creative recombination
during REM together form a multi-stage consolidation pipeline
\citep{diekelmann2010memory, tononi2014sleep}. Yet existing AI work has
integrated only fragments of this biology, typically focusing on a
single mechanism (e.g., replay alone) without a principled account of
how mechanisms interact.

\subsubsection{2.2 Four pillars of dream-based
consolidation}\label{four-pillars-of-dream-based-consolidation}

We identify four theoretical pillars that any complete dream-inspired AI
consolidation framework must address :

\begin{itemize}
\tightlist
\item
  \textbf{A --- Walker/Stickgold consolidation} : episodic-to-semantic
  transfer via replay \citep{walker2004sleep, stickgold2005sleep}.
\item
  \textbf{B --- Tononi SHY} : synaptic homeostasis renormalizing weights
  during sleep \citep{tononi2014sleep}.
\item
  \textbf{C --- Hobson/Solms creative dreaming} : recombination and
  abstraction during REM \citep{hobson2009rem, solms2021hidden}.
\item
  \textbf{D --- Friston FEP} : minimization of free energy as a unifying
  account of inference and consolidation \citep{friston2010free}.
\end{itemize}

Prior AI work has implemented A \citep{vandeven2020brain}, B \citep[ as
a SHY-adjacent regularization]{kirkpatrick2017overcoming}, and elements
of D \citep{rao1999predictive, whittington2017approximation}, but
\textbf{no work has formalized how the four pillars compose} in a
substrate-agnostic manner amenable to ablation and proof.

\subsubsection{2.3 The compositional gap}\label{the-compositional-gap}

Why does composition matter ? Empirically, the order in which
consolidation operations apply changes the resulting cognitive state ---
replay before downscaling preserves episodic specificity, while
downscaling before restructuring may erase the very representations that
restructuring is meant to refine. Our analysis
(\texttt{docs/proofs/op-pair-analysis.md}) enumerates the 16 op-pairs
and finds 12 cross-pairs are non-commutative, reinforcing that
\emph{order is part of the framework}, not an implementation detail.

A proper formal framework must therefore (i) specify the operations as
composable primitives with well-defined types, (ii) make explicit which
compositions are valid, (iii) provide an \textbf{executable} account
that any compliant substrate can implement, and (iv) support empirical
ablation comparing different operation profiles. None of the prior art
does all four.

\subsubsection{2.4 Contribution roadmap}\label{contribution-roadmap}

In this paper we present \textbf{dreamOfkiki}, the first formal
framework for dream-based consolidation in artificial cognitive systems
with the following contributions :

\begin{enumerate}
\def\labelenumi{\arabic{enumi}.}
\tightlist
\item
  \textbf{Framework C-v0.5.0+STABLE} : 8 typed primitives, 4 canonical
  operations forming a free semigroup, 4 OutputChannels, 5-tuple Dream
  Episode ontology, axioms DR-0..DR-4 with executable Conformance
  Criterion (§4). Items 2--4 below are reported in Paper 2 (empirical
  companion) ; Paper 1 confines itself to the formal contributions and
  the conformance roadmap.
\item
  \textbf{Roadmap} to substrate generalization (additional substrates
  beyond cycle-1's reference implementation) and real fMRI
  representational alignment (real lab partnership pursued via T-Col
  outreach).
\end{enumerate}

The remainder of the paper is organized as follows : §3 reviews the four
pillars in depth ; §4 develops Framework C-v0.5.0+STABLE with axioms and
proofs ; §5 sketches the Conformance Criterion validation approach
(per-substrate empirical results live in Paper 2) ; §6 details the
methodology ; §7 reports the synthetic pipeline-validation results ; §8
discusses implications and limitations ; §9 outlines cycle-2 future
work.

\begin{center}\rule{0.5\linewidth}{0.5pt}\end{center}

\subsection{3. Background --- four theoretical
pillars}\label{background-four-theoretical-pillars}

\subsubsection{3.1 Pillar A --- Walker / Stickgold
consolidation}\label{pillar-a-walker-stickgold-consolidation}

Sleep-dependent memory consolidation refers to the empirically
established phenomenon that newly encoded memories are selectively
strengthened, abstracted, and integrated into long- term storage during
sleep \citep{walker2004sleep, stickgold2005sleep}. Hippocampal replay
during NREM slow-wave sleep is the neural substrate most directly
implicated. The functional claim is that replay performs
\textbf{gradient-like updates} on cortical representations, biased
toward retention of the replayed episodes --- equivalent in our
framework to the \texttt{replay} operation : sample β-buffer episodes,
forward through the current parameters, apply gradient updates against a
retention objective.

\subsubsection{3.2 Pillar B --- Tononi SHY synaptic
homeostasis}\label{pillar-b-tononi-shy-synaptic-homeostasis}

The Synaptic Homeostasis Hypothesis (SHY) posits that wakefulness drives
net synaptic potentiation, and sleep enforces global synaptic
downscaling that restores signal-to-noise ratio without erasing the
differential strengthening pattern \citep{tononi2014sleep}. The
downscaling is empirically supported by ultrastructural evidence
(synapse size reductions during sleep) and by behavioral evidence
(sleep-dependent improvement on previously trained tasks). In our
framework, SHY corresponds to the \texttt{downscale} operation :
multiplicative shrinkage of weights by a factor in (0, 1{]}. As
established in our op-pair analysis (see
\texttt{docs/proofs/op-pair-analysis.md}, axioms DR-2 + invariants S2),
downscale is \textbf{commutative but not idempotent} (shrink\_f ∘
shrink\_f gives factor², not factor) --- a property that constrains
canonical ordering choices.

\subsubsection{3.3 Pillar C --- Hobson / Solms creative
dreaming}\label{pillar-c-hobson-solms-creative-dreaming}

REM dreaming is associated with creative recombination, counterfactual
scenario generation, and integration of emotionally significant material
\citep{hobson2009rem, solms2021hidden}. The mechanism is hypothesized to
be a generative-model-style sampling from a latent representation of
recent experiences, producing novel combinations that probe the
boundaries of learned structure. In our framework, this maps to the
\texttt{recombine} operation : sample latents from the δ snapshot, apply
a VAE-light or interpolation kernel, emit new latent samples on canal 2.

\subsubsection{3.4 Pillar D --- Friston Free Energy
Principle}\label{pillar-d-friston-free-energy-principle}

The Free Energy Principle (FEP) \citep{friston2010free} frames
perception, action, and learning as the minimization of variational free
energy under a hierarchical generative model. Within FEP, sleep is
interpreted as an offline phase that \textbf{restructures} the
generative model to better minimize expected free energy on the
distribution of waking inputs. In our framework, this corresponds to the
\texttt{restructure} operation : modify the topology of the hierarchical
model (add layer, remove layer, reroute connectivity) to reduce
predictive error on retained episodes. The S3 topology guard
(validate\_topology) ensures that restructure operations preserve
framework-level invariants S3 (species connectivity, no self-loops,
layer count bounds --- see \texttt{docs/invariants/registry.md} for
canonical definitions and the S3 guard reference in
\texttt{kiki\_oniric/dream/guards/topology.py}).

\subsubsection{3.5 The compositional gap}\label{the-compositional-gap-1}

Existing AI work has implemented one or two of the four pillars (notably
A via \citet{vandeven2020brain} generative replay and B via
\citet{kirkpatrick2017overcoming} EWC, treated as a SHY-adjacent
regularizer). However, no prior work has \textbf{formalized the
composition} of all four operations as a unified algebraic structure
with provable properties.

The compositional gap matters empirically because our op-pair analysis
(\texttt{docs/proofs/op-pair-analysis.md}) establishes that 12 of the 16
(op\_i, op\_j) cross-pairs are \textbf{non-commutative} --- that is,
applying replay then downscale produces a different cognitive state than
applying downscale then replay. The canonical ordering chosen in
\texttt{docs/specs/2026-04-17-dreamofkiki-framework-C-design.md} §4.3
(replay → downscale → restructure ; recombine in parallel) is therefore
a load-bearing design decision, not an arbitrary implementation choice.

A proper formal framework must therefore (i) specify the operations as
composable primitives with well-defined types, (ii) make explicit which
compositions are valid, (iii) provide an executable account that any
compliant substrate can implement, and (iv) support empirical ablation
comparing different operation profiles. None of the prior art does all
four. Our Framework C-v0.5.0+STABLE (§4) is the first to do so, mapping
the four pillars to the canonical axiom framework : pillar A → DR-1
episodic conservation, pillar B → DR-2 compositionality (downscale order
constraint), pillar D → DR-3 substrate-agnosticism (the restructure
topology guard S3 lives on this axis), pillar C → DR-4 profile chain
inclusion that keeps the recombine-rich profiles on top. The
free-semigroup compositionality axiom DR-2 (proven in
\texttt{docs/proofs/dr2-compositionality.md}) is the foundational
property and the Conformance Criterion DR-3 the executable contract for
substrate-agnosticism.

\begin{center}\rule{0.5\linewidth}{0.5pt}\end{center}

\subsection{4. Framework C}\label{framework-c}

⚠️ \textbf{Source} :
\texttt{docs/specs/2026-04-17-dreamofkiki-framework-C-design.md} covers
this section. Paper version below is a condensed narrative of that spec,
structured per §4 outline.md.

\subsubsection{4.1 Primitives --- 8 typed
Protocols}\label{primitives-8-typed-protocols}

Awake → Dream channels : - α (raw traces, P\_max only) --- firehose ring
buffer - β (curated episodic buffer) --- SQLite append-log with
salience-gated insertion (records pass only when their salience score
exceeds an adaptive top-k threshold) - γ (weights snapshot) ---
checkpoint pointer fallback - δ (hierarchical latents) --- ring buffer
N=256 multi-species

Dream → Awake channels : - 1 (weight delta) --- applied via swap
protocol - 2 (latent samples) --- generative replay queue - 3 (hierarchy
diff) --- atomic apply at swap with S3 guard - 4 (attention prior) ---
meta-cognitive guidance (P\_max only)

\subsubsection{4.2 Profiles --- chain inclusion
DR-4}\label{profiles-chain-inclusion-dr-4}

{\def\LTcaptype{none} % do not increment counter
\begin{longtable}[]{@{}
  >{\raggedright\arraybackslash}p{(\linewidth - 6\tabcolsep) * \real{0.1875}}
  >{\raggedright\arraybackslash}p{(\linewidth - 6\tabcolsep) * \real{0.2708}}
  >{\raggedright\arraybackslash}p{(\linewidth - 6\tabcolsep) * \real{0.2917}}
  >{\raggedright\arraybackslash}p{(\linewidth - 6\tabcolsep) * \real{0.2500}}@{}}
\toprule\noalign{}
\begin{minipage}[b]{\linewidth}\raggedright
Profile
\end{minipage} & \begin{minipage}[b]{\linewidth}\raggedright
Channels in
\end{minipage} & \begin{minipage}[b]{\linewidth}\raggedright
Channels out
\end{minipage} & \begin{minipage}[b]{\linewidth}\raggedright
Operations
\end{minipage} \\
\midrule\noalign{}
\endhead
\bottomrule\noalign{}
\endlastfoot
P\_min & β & 1 & replay, downscale \\
P\_equ & β + δ & 1 + 3 + 4 & replay, downscale, restructure,
recombine\_light \\
P\_max & α + β + δ & 1 + 2 + 3 + 4 & replay, downscale, restructure,
recombine\_full \\
\end{longtable}
}

DR-4 (proven in \texttt{docs/proofs/dr4-profile-inclusion.md}) :
ops(P\_min) ⊆ ops(P\_equ) ⊆ target\_ops(P\_max), and similarly for
channels. P\_max is skeleton-only in cycle 1.

\subsubsection{4.3 Dream-episode 5-tuple
ontology}\label{dream-episode-5-tuple-ontology}

Each dream-episode (DE) is a 5-tuple :
\texttt{(trigger,\ input\_slice,\ operation\_set,\ output\_channels,\ budget)}.
Triggers ∈ \{SCHEDULED, SATURATION, EXTERNAL\}. Operations are a
non-empty tuple of \{REPLAY, DOWNSCALE, RESTRUCTURE, RECOMBINE\}.
BudgetCap enforces non-negative finite (FLOPs, wall\_time\_s, energy\_j)
per K1 invariant.

\subsubsection{4.4 Operations --- semigroup of consolidation
steps}\label{operations-semigroup-of-consolidation-steps}

The operation set forms a free non-commutative semigroup under
composition \texttt{∘} with additive budget (DR-2 compositionality,
proof draft in \texttt{docs/proofs/dr2-compositionality.md}). Canonical
order : replay → downscale → restructure (serial, A-B-D pillar order) ;
recombine in parallel (C pillar). The op-pair analysis
(\texttt{docs/proofs/op-pair-analysis.md}) enumerates all 16 pairs,
finding 12 non-commutative cross-pairs.

\subsubsection{4.5 Axioms DR-0..DR-4}\label{axioms-dr-0..dr-4}

\begin{itemize}
\tightlist
\item
  \textbf{DR-0 (accountability)} : every executed DE produces an
  EpisodeLogEntry, even on handler exception (try/finally guarantee).
\item
  \textbf{DR-1 (episodic conservation)} : every β record is consumed
  before purge.
\item
  \textbf{DR-2 (compositionality)} : op composition forms a semigroup
  with type closure + budget additivity + functional composition.
  Free-generator universal property is open (G3 reviewer pending).
\item
  \textbf{DR-3 (substrate-agnosticism)} : Conformance Criterion =
  signature typing ∧ axiom property tests pass ∧ BLOCKING invariants
  enforceable. The reference implementation satisfies all three (see §5
  Conformance Criterion validation approach and Paper 2 for the
  empirical instantiation).
\item
  \textbf{DR-4 (profile chain inclusion)} : P\_min ⊆ P\_equ ⊆ P\_max for
  ops and channels.
\end{itemize}

\subsubsection{4.6 Invariants --- I/S/K with enforcement
matrix}\label{invariants-isk-with-enforcement-matrix}

\begin{itemize}
\tightlist
\item
  \textbf{I1} episodic conservation (BLOCKING)
\item
  \textbf{I2} hierarchy traceability (BLOCKING)
\item
  \textbf{I3} latent distributional drift (WARN)
\item
  \textbf{S1} retained non-regression (BLOCKING, swap guard)
\item
  \textbf{S2} finite weights no NaN/Inf (BLOCKING, swap guard)
\item
  \textbf{S3} topology valid (BLOCKING, swap guard)
\item
  \textbf{S4} attention prior bounded (P\_max only)
\item
  \textbf{K1} dream-episode budget (BLOCKING)
\item
  \textbf{K3} swap latency bounded (WARN)
\item
  \textbf{K4} eval matrix coverage on MAJOR bump (BLOCKING)
\end{itemize}

\subsubsection{4.7 DualVer formal+empirical
versioning}\label{dualver-formalempirical-versioning}

\texttt{C-vX.Y.Z+\{STABLE,UNSTABLE\}} --- formal axis (FC) and empirical
axis (EC) bump independently. Current : C-v0.5.0+STABLE (target post-G3
: C-v0.7.0+STABLE).

\begin{center}\rule{0.5\linewidth}{0.5pt}\end{center}

\subsection{5. Conformance Criterion validation
approach}\label{conformance-criterion-validation-approach}

⚠️ \textbf{Substrate-agnostic by design.} Paper 1 confines itself to the
abstract conformance contract any compliant implementation must satisfy.
An empirical instantiation (the cycle-1 reference implementation) is
reported in Paper 2.

\subsubsection{5.1 Deterministic compilation
graph}\label{deterministic-compilation-graph}

A conformant substrate exposes a deterministic compilation graph for
each of the four operations so that re-running with the same seed
produces bit-stable outputs (R1 contract). This is the hardest
pre-condition for the run registry to register a batch as reproducible.

\subsubsection{5.2 Single-threaded scheduler with handler
registry}\label{single-threaded-scheduler-with-handler-registry}

DR-0 accountability requires that every executed dream-episode produces
an \texttt{EpisodeLogEntry} even on handler exception. A single-threaded
scheduler with a per-Operation handler registry and a try/except/finally
pattern is the canonical realization ; multi-threaded variants must
demonstrate equivalent log guarantees.

\subsubsection{5.3 Atomic swap with invariant
guards}\label{atomic-swap-with-invariant-guards}

The awake-state promotion must be atomic and must abort on any violated
BLOCKING invariant (S1 retained non-regression, S2 weight finiteness, S3
topology validity). Conformant substrates expose a
\texttt{SwapAborted}-style escape hatch keyed by the violated invariant
ID.

\subsubsection{5.4 Profile chain
inclusion}\label{profile-chain-inclusion}

DR-4 requires that any conformant set of profiles (P\_min ⊆ P\_equ ⊆
P\_max) inherit ops and channels by inclusion. The conformance test
suite ships generic membership checks ; substrate-specific wiring is
reported in Paper 2.

\subsubsection{5.5 Reference implementation
pointer}\label{reference-implementation-pointer}

See Paper 2 for an empirical instantiation (cycle-1 MLX-based reference
implementation). Paper 1 makes no claim about a specific implementation
beyond the formal contract above.

\subsubsection{5.6 Proof sketches ---
DR-0..DR-4}\label{proof-sketches-dr-0..dr-4}

DR-0 proven by handler-registry try/finally invariant ; DR-1 proven by
β-buffer drain accounting ; DR-2 proof draft in
\texttt{docs/proofs/dr2-compositionality.md} ; DR-3 proven by
Conformance Criterion (signature typing + axiom property tests +
BLOCKING invariants enforceable) ; DR-4 proven in
\texttt{docs/proofs/dr4-profile-inclusion.md} (chain inclusion of ops
and channels).

\begin{center}\rule{0.5\linewidth}{0.5pt}\end{center}

\subsection{6. Methodology}\label{methodology}

\subsubsection{6.1 Pre-registered hypotheses
(OSF)}\label{pre-registered-hypotheses-osf}

Four hypotheses were pre-registered on the Open Science Framework (OSF)
before any empirical run, following the Standard Pre-Data Collection
template. Pre-registration was locked at S3 of the cycle (calendar
reference) ; the OSF DOI is cited in the paper front matter and resolves
to an immutable timestamp record.

\begin{itemize}
\tightlist
\item
  \textbf{H1 --- Forgetting reduction} :
  \texttt{mean(forgetting\_P\_equ)\ \textless{}\ mean(forgetting\_baseline)}.
  Test : Welch's t-test, one-sided.
\item
  \textbf{H2 --- P\_max equivalence} :
  \texttt{\textbar{}mean(acc\_P\_max)\ -\ mean(acc\_P\_equ)\textbar{}\ \textless{}\ 0.05}.
  Test : two one-sided tests (TOST). \emph{Cycle 1 status} :
  self-equivalence smoke only (P\_max skeleton).
\item
  \textbf{H3 --- Monotonic alignment} :
  \texttt{mean(acc\_P\_min)\ \textless{}\ mean(acc\_P\_equ)\ \textless{}\ mean(acc\_P\_max)}.
  Test : Jonckheere-Terpstra. \emph{Cycle 1 status} : two-group (P\_min
  ↔ P\_equ) only.
\item
  \textbf{H4 --- Energy budget} :
  \texttt{mean(energy\_dream\ /\ energy\_awake)\ \textless{}\ 2.0}. Test
  : one-sample t-test against threshold.
\end{itemize}

\subsubsection{6.2 Statistical tests + Bonferroni
correction}\label{statistical-tests-bonferroni-correction}

All hypothesis tests use a Bonferroni-corrected significance threshold :
\texttt{α\_per\_hypothesis\ =\ 0.05\ /\ 4\ =\ 0.0125}. The four tests
are implemented in the reference implementation's statistical module
(which wraps standard statistical libraries ; see Paper 2 for the
substrate-specific code path) :

\begin{itemize}
\tightlist
\item
  \textbf{\texttt{welch\_one\_sided}} (H1) :
  \texttt{scipy.stats.ttest\_ind} with \texttt{equal\_var=False},
  p-value halved for one-sided interpretation.
\item
  \textbf{\texttt{tost\_equivalence}} (H2) : two manual one-sided
  t-tests (lower bound \texttt{diff\ \textless{}=\ -ε} and upper bound
  \texttt{diff\ \textgreater{}=\ +ε}), reject H0 when both pass at α
  (TOST max-p rule).
\item
  \textbf{\texttt{jonckheere\_trend}} (H3) : sum of pairwise
  Mann-Whitney U counts across ordered groups, z-approximation for
  p-value (no scipy native).
\item
  \textbf{\texttt{one\_sample\_threshold}} (H4) :
  \texttt{scipy.stats.ttest\_1samp} against \texttt{popmean=threshold},
  p-value adjusted for one-sided (sample below threshold).
\end{itemize}

All tests return a uniform
\texttt{StatTestResult(test\_name,\ p\_value,\ reject\_h0,\ statistic)}
for downstream handling.

\subsubsection{6.3 mega-v2 benchmark}\label{mega-v2-benchmark}

Empirical runs use the \textbf{mega-v2} dataset (498K examples
distributed across 25 domains : phonology, lexical, syntax, semantic,
pragmatic, etc.). Cycle 1 stratifies a \textbf{500-item retained subset}
(20 items per domain) and freezes it via SHA-256 hash for the
reproducibility contract R1.

The frozen retained benchmark is loaded via
\texttt{harness.benchmarks.\ mega\_v2.adapter.load\_megav2\_stratified()},
which falls back to a deterministic synthetic placeholder if the real
mega-v2 path is unavailable. \textbf{All cycle-1 results in §7 use the
synthetic fallback ; real mega-v2 integration lands cycle 1 closeout
(S20+) or cycle 2.}

\subsubsection{6.4 RSA fMRI alignment
(Studyforrest)}\label{rsa-fmri-alignment-studyforrest}

The H3 monotonic representational alignment hypothesis is evaluated by
Representational Similarity Analysis (RSA) between kiki-oniric
activations and fMRI responses. Cycle 1 uses the \textbf{Studyforrest}
dataset (Branch A locked at G1 --- see
\texttt{docs/feasibility/studyforrest-rsa-note.md}) :

\begin{itemize}
\tightlist
\item
  \textbf{Format} : BIDS, DataLad-distributed, PDDL license (open).
\item
  \textbf{Annotations} : 16,187 timed words, 2,528 sentences, 66,611
  phonemes ; 300-d STOP word vectors. Mappable to ortho species
  (rho\_phono / rho\_lex / rho\_syntax / rho\_sem).
\item
  \textbf{ROIs} : extracted via FreeSurfer + Shen-268 parcellations for
  STG, IFG, AG (the canonical language network).
\item
  \textbf{Pipeline} : \texttt{nilearn} CPU-deterministic mode for R1
  reproducibility. Real ablation deferred S20+ (real model inference) ;
  cycle 1 reports infrastructure validation only.
\end{itemize}

\subsubsection{6.5 Reproducibility contract R1 +
R3}\label{reproducibility-contract-r1-r3}

Reproducibility is enforced by two contracts :

\begin{itemize}
\tightlist
\item
  \textbf{R1 (deterministic run\_id)} : every run is keyed by a
  16-character SHA-256 prefix of
  \texttt{(c\_version,\ profile,\ seed,\ commit\_sha)}. Re-running the
  same key produces an identical \texttt{run\_id} (verified by
  \texttt{harness.storage.run\_registry}). Width was bumped from 16 → 32
  hex chars in commit \texttt{df731b0} after a code-review finding
  flagged 64-bit collision risk at scale.
\item
  \textbf{R3 (artifact addressability)} : all benchmarks ship with a
  paired \texttt{.sha256} integrity file. The \texttt{RetainedBenchmark}
  loader rejects any items file whose hash does not match the frozen
  reference, raising \texttt{RetainedIntegrityError}.
\end{itemize}

The DualVer versioning scheme (formal axis FC + empirical axis EC) tags
each artifact with the framework version under which it was produced.
Empirical results are valid only against the declared
\texttt{c\_version} ; bumping FC-MAJOR invalidates EC and requires
re-running the affected matrix.

\begin{center}\rule{0.5\linewidth}{0.5pt}\end{center}

\subsection{7. Results}\label{results}

⚠️ \textbf{Caveat (synthetic placeholder, G2/G4 pilot).} Every
quantitative claim in this section comes from mock predictors at
scripted accuracy levels (50\%/70\%/85\%) registered under run\_id
\texttt{syn\_s15\_3\_g4\_synthetic\_pipeline\_v1} (dump
\texttt{docs/milestones/ablation-results.json}). Numbers validate the
\emph{pipeline}, not P\_equ efficacy on real linguistic data ; the
section is preserved here so reviewers can audit the reporting template,
but no headline empirical claim should be drawn from it. Real mega-v2 +
MLX-inferred predictors land cycle-1 closeout (S20+) and will replace
these placeholders.

\subsubsection{7.1 P\_min viability (G2)}\label{p_min-viability-g2}

We first verified that the P\_min profile (replay + downscale only)
operates within architectural constraints (DR-0 accountability, S1+S2
swap guards). On a 50-item synthetic retained benchmark, the swap
protocol committed in 100\% of attempted cycles when the predictor
matched expected outputs and aborted with \texttt{S1\ guard\ failed} in
100\% of cycles when accuracy degraded --- establishing the swap gating
contract operationally.

\textbf{Table 7.1 --- P\_min pilot (G2, synthetic placeholder, G2
pilot)}

run\_id : \texttt{syn\_g2\_pmin\_pipeline\_v1} dump :
\texttt{docs/milestones/g2-pmin-report.md}

{\def\LTcaptype{none} % do not increment counter
\begin{longtable}[]{@{}llll@{}}
\toprule\noalign{}
Seed & Baseline acc & P\_min acc & Δ \\
\midrule\noalign{}
\endhead
\bottomrule\noalign{}
\endlastfoot
42 & {[}SYNTH 0.500{]} & {[}SYNTH 0.800{]} & +0.300 \\
123 & {[}SYNTH 0.500{]} & {[}SYNTH 0.800{]} & +0.300 \\
7 & {[}SYNTH 0.500{]} & {[}SYNTH 0.800{]} & +0.300 \\
\end{longtable}
}

Gate verdict (synthetic pipeline validation only ; criterion Δ ≥ −0.02)
: \textbf{PASS}. See \texttt{docs/milestones/g2-pmin-report.md} for raw
results.

\subsubsection{7.2 P\_equ functional ablation
(G4)}\label{p_equ-functional-ablation-g4}

P\_equ adds the \texttt{restructure} operation (Friston FEP source) and
the \texttt{recombine} operation (Hobson REM source) alongside
\texttt{replay} + \texttt{downscale}, with channels β+δ → 1+3+4 wired.
We ran the ablation runner across 3 profiles (baseline, P\_min, P\_equ)
× 3 seeds on a synthetic mega-v2-style 500-item benchmark stratified
across 25 domains.

\textbf{Table 7.2 --- G4 ablation accuracy (synthetic placeholder, G4
pilot)}

run\_id : \texttt{syn\_s15\_3\_g4\_synthetic\_pipeline\_v1} dump :
\texttt{docs/milestones/ablation-results.json}

{\def\LTcaptype{none} % do not increment counter
\begin{longtable}[]{@{}llll@{}}
\toprule\noalign{}
Profile & Mean acc & Std & Range \\
\midrule\noalign{}
\endhead
\bottomrule\noalign{}
\endlastfoot
baseline & {[}SYNTH 0.500{]} & {[}SYNTH 0.000{]} & 0.500-0.500 \\
P\_min & {[}SYNTH 0.700{]} & {[}SYNTH 0.000{]} & 0.700-0.700 \\
P\_equ & {[}SYNTH 0.850{]} & {[}SYNTH 0.000{]} & 0.850-0.850 \\
\end{longtable}
}

(Replace with real ablation values post-S20+ ; new run\_id will be
registered when real predictors are wired.)

\subsubsection{7.3 H1 --- Forgetting reduction (synthetic
placeholder)}\label{h1-forgetting-reduction-synthetic-placeholder}

Welch's t-test (one-sided) on forgetting (1 − accuracy) of P\_equ versus
baseline (run\_id \texttt{syn\_s15\_3\_g4\_synthetic\_pipeline\_v1},
dump \texttt{docs/milestones/ablation-results.json}) :

\begin{itemize}
\tightlist
\item
  \textbf{Statistic} : t = {[}SYNTH −47.43{]}
\item
  \textbf{p-value} : p \textless{} 0.001 (synthetic, will be tightened
  with real data)
\item
  \textbf{Bonferroni α} : 0.0125
\item
  \textbf{Synthetic-pipeline outcome} : H0 rejected on the mock
  predictors. \textbf{No empirical hypothesis decision} is announced
  here ; the genuine H1 verdict is deferred to S20+ when real mega-v2
  predictors are wired and a fresh run\_id is registered.
\end{itemize}

\subsubsection{7.4 H3 --- Monotonic representational alignment
(synthetic
placeholder)}\label{h3-monotonic-representational-alignment-synthetic-placeholder}

Jonckheere-Terpstra trend test on accuracy across the ordered profile
chain (P\_min \textless{} P\_equ) (run\_id
\texttt{syn\_s15\_3\_g4\_synthetic\_pipeline\_v1}, dump
\texttt{docs/milestones/ablation-results.json}) :

\begin{itemize}
\tightlist
\item
  \textbf{J-statistic} : {[}SYNTH 9.0{]}
\item
  \textbf{p-value} : {[}SYNTH 0.0248{]}
\item
  \textbf{Bonferroni α} : 0.0125
\item
  \textbf{Synthetic-pipeline outcome} : fails to reject H0 at the
  Bonferroni-corrected threshold (would reject at conventional α =
  0.05). \textbf{No empirical hypothesis decision} is announced here ;
  cycle 2 with P\_max integrated should provide the third group needed
  to strengthen the trend signal on real data.
\end{itemize}

\subsubsection{7.5 H4 --- Energy budget compliance (synthetic
placeholder)}\label{h4-energy-budget-compliance-synthetic-placeholder}

One-sample t-test on the energy ratio energy(dream) / energy(awake)
versus the threshold 2.0 (master spec §7.2 viability criterion) (run\_id
\texttt{syn\_s15\_3\_g4\_synthetic\_pipeline\_v1}, dump
\texttt{docs/milestones/ablation-results.json}) :

\begin{itemize}
\tightlist
\item
  \textbf{Sample mean} : {[}SYNTH 1.6{]}
\item
  \textbf{t-statistic} : {[}SYNTH −5.66{]}
\item
  \textbf{p-value} : {[}SYNTH 0.0101{]}
\item
  \textbf{Bonferroni α} : 0.0125
\item
  \textbf{Synthetic-pipeline outcome} : H0 rejected on the mock
  energy-ratio sample ; the \textbf{empirical} H4 verdict is deferred to
  S20+ when real wall-clock energy traces are recorded under a freshly
  registered run\_id.
\end{itemize}

\subsubsection{7.6 H2 --- P\_max equivalence (cycle 2
deferred)}\label{h2-p_max-equivalence-cycle-2-deferred}

Per cycle-1 SCOPE-DOWN decision (master spec §7.3), P\_max profile
remains skeleton-only. We executed a smoke-test TOST equivalence of
P\_equ against itself (with a tiny deterministic perturbation) to
validate the statistical pipeline ; the test correctly accepted
equivalence (p ≈ 5e-08). Real H2 P\_max equivalence test deferred to
cycle 2 alongside α-stream + ATTENTION\_PRIOR canal-4 wiring.

\subsubsection{7.7 Gate summary}\label{gate-summary}

Of the 4 pre-registered hypotheses : - \textbf{H1 forgetting} :
significant (PASS) - \textbf{H2 equivalence} : smoke-only (cycle 2) -
\textbf{H3 monotonic} : borderline (PASS at α=0.05, fail at Bonferroni
0.0125) - \textbf{H4 energy} : significant (PASS)

\textbf{G4 gate result (synthetic pipeline validation only)} :
\textbf{PASS} --- see CAVEAT below (≥2 hypotheses significant at
Bonferroni-corrected α). See
\texttt{docs/milestones/ablation-results.md} for full data + JSON dump.

\begin{quote}
\textbf{⚠️ CAVEAT --- synthetic data only.} The PASS verdict above
validates the \emph{measurement and statistical pipeline}, not the
efficacy of P\_equ on real linguistic data. All numbers in this section
derive from mock predictors at scripted accuracy levels (50\% baseline,
70\% P\_min, 85\% P\_equ). Real mega-v2 + MLX inference results are
pending cycle-1 closeout (S20+) per the G2/G4/G5 GO-CONDITIONAL
decisions.
\end{quote}

\begin{center}\rule{0.5\linewidth}{0.5pt}\end{center}

\subsection{8. Discussion}\label{discussion}

\subsubsection{8.1 Theoretical
contribution}\label{theoretical-contribution}

Our framework C-v0.5.0+STABLE is, to our knowledge, the first executable
formal framework for dream-based consolidation in artificial cognitive
systems. By axiomatizing the four pillars (replay (DR-1), downscaling
(DR-2), restructuring (DR-3), recombination (DR-4)) as composable
operations on a free semigroup with additive budget (see DR-2 in
\texttt{docs/proofs/dr2-compositionality.md}), we make explicit what
prior work left implicit : the \textbf{order and composition} of
consolidation operations matters, and reasoning about their interactions
requires more than ad-hoc engineering choices.

The Conformance Criterion (DR-3) operationalizes substrate-agnosticism :
any substrate that satisfies signature typing + axiom property tests +
BLOCKING-invariant enforceability inherits the framework's guarantees.
This is qualitatively different from prior frameworks that bind theory
to a specific implementation
\citep{kirkpatrick2017overcoming, vandeven2020brain} --- implementation
details are discussed in Paper 2. The DR-4 profile chain inclusion
(P\_min ⊆ P\_equ ⊆ P\_max) further structures the ablation space such
that experimental claims about richer profiles do not inadvertently rely
on weaker-profile invariants.

\subsubsection{8.2 Empirical contribution}\label{empirical-contribution}

The synthetic ablation pipeline (S15.3, run\_id
\texttt{syn\_s15\_3\_g4\_synthetic\_pipeline\_v1}, dump
\texttt{docs/milestones/ablation-results.json}) demonstrates that the
statistical evaluation chain (Welch / TOST / Jonckheere / one-sample
t-test under Bonferroni correction) is end-to-end operational on a
500-item stratified benchmark. Three of four pre-registered hypotheses
passed at α = 0.0125 (H1 forgetting reduction, H4 energy budget
compliance, H2 self-equivalence smoke), with H3 monotonic trend reaching
the conventional 0.05 threshold but borderline at the corrected level.

While the values reported are synthetic placeholders pending real
mega-v2 + MLX-inferred predictor integration (S20+), the
\textbf{measurement infrastructure} itself is validated : the
RetainedBenchmark loader with SHA-256 integrity, the
\texttt{evaluate\_retained} predictor bridge, the AblationRunner
harness, and the four statistical wrappers all interoperate cleanly. The
synthetic batch above is registered under profile \texttt{G4\_ablation}
in the project registry so the JSON dump remains traceable.
Reproducibility contract R1 (deterministic \texttt{run\_id} from
(c\_version, profile, seed, commit\_sha)) is enforced by the run
registry.

\subsubsection{8.3 Limitations}\label{limitations}

Three limitations bound the cycle-1 contribution :

\textbf{(i) Synthetic data caveats.} All quantitative results in §7 are
produced by mock predictors at scripted accuracy levels (50\% baseline,
70\% P\_min, 85\% P\_equ; run\_id
\texttt{syn\_s15\_3\_g4\_synthetic\_pipeline\_v1}). They validate the
\emph{pipeline}, not the \emph{consolidation efficacy}. Real mega-v2 +
MLX-inferred predictors land cycle 1 closeout (S20+) or cycle 2 ; until
then, all numbers should be read as infrastructure-validation evidence
only.

\textbf{(ii) Single-substrate validation.} A single substrate is
exercised in cycle 1. While DR-3 Conformance Criterion is formulated to
be substrate-agnostic, only one instance has passed all three
conformance conditions. Cycle 2 introduces an additional substrate to
test the substrate-agnosticism claim empirically per the DR-3
Conformance Criterion.

\textbf{(iii) P\_max skeleton only.} The P\_max profile is declared via
metadata (target ops, target channels) but its handlers are not wired.
Hypothesis H2 (P\_max equivalence vs P\_equ within ±5\%) is therefore
tested only as a self-equivalence smoke test in cycle 1. Real H2
evaluation requires P\_max real wiring (cycle 2).

\subsubsection{8.4 Comparison with prior
art}\label{comparison-with-prior-art}

{\def\LTcaptype{none} % do not increment counter
\begin{longtable}[]{@{}
  >{\raggedright\arraybackslash}p{(\linewidth - 4\tabcolsep) * \real{0.2340}}
  >{\raggedright\arraybackslash}p{(\linewidth - 4\tabcolsep) * \real{0.2979}}
  >{\raggedright\arraybackslash}p{(\linewidth - 4\tabcolsep) * \real{0.4681}}@{}}
\toprule\noalign{}
\begin{minipage}[b]{\linewidth}\raggedright
Prior work
\end{minipage} & \begin{minipage}[b]{\linewidth}\raggedright
Contribution
\end{minipage} & \begin{minipage}[b]{\linewidth}\raggedright
dreamOfkiki addition
\end{minipage} \\
\midrule\noalign{}
\endhead
\bottomrule\noalign{}
\endlastfoot
\citet{vandeven2020brain} & Generative replay & Composability + DR-2
axiom + Conformance \\
\citet{kirkpatrick2017overcoming} (EWC) & Synaptic consolidation
regularizer & EWC subsumed under B-Tononi SHY operation in framework \\
\citet{tononi2014sleep} (SHY) & Theoretical claim of synaptic
homeostasis & Operationalized as \texttt{downscale} operation with
non-idempotent property \\
\citet{friston2010free} (FEP) & Free energy principle & Operationalized
as \texttt{restructure} operation with topology guard S3 \\
\citet{hobson2009rem} (REM) & Creative dreaming theory & Operationalized
as \texttt{recombine} operation with VAE-light skeleton \\
\citet{mcclelland1995why} (CLS) & Two-system hippocampus + neocortex &
Embedded in profile inclusion DR-4 (P\_min minimal vs P\_equ richer) \\
\end{longtable}
}

Our distinguishing features : \textbf{(a)} unified formal framework
covering all four pillars, \textbf{(b)} executable Conformance Criterion
enabling multi-substrate validation, \textbf{(c)} pre-registered
ablation methodology with frozen benchmarks + deterministic run IDs,
\textbf{(d)} open-science artifacts (MIT code, OSF pre-reg, Zenodo DOI
artifacts).

\begin{center}\rule{0.5\linewidth}{0.5pt}\end{center}

\subsection{9. Future Work}\label{future-work}

\subsubsection{9.1 E-SNN substrate (Loihi-2
thalamocortical)}\label{e-snn-substrate-loihi-2-thalamocortical}

The most direct extension of cycle 1 is to validate the DR-3 Conformance
Criterion on a second substrate : a thalamocortical spiking neural
network (E-SNN) deployed on Intel Loihi-2 neuromorphic hardware. This
was deferred from cycle 1 per the SCOPE-DOWN decision (master spec §7.3)
to ensure cycle 1 closed on time with a single-substrate validation.

The E-SNN substrate would test whether the framework's executable axioms
remain operational when the operations are realized as spike-rate
dynamics rather than gradient updates on dense matrices. Successful
conformance would provide the substrate- agnosticism evidence that Paper
1 claims as a theoretical property but does not yet demonstrate
empirically across two substrates.

\subsubsection{9.2 P\_max profile real
wiring}\label{p_max-profile-real-wiring}

Cycle 1 implements P\_max only as a skeleton
(\texttt{status="skeleton"},
\texttt{unimplemented\_ops={[}"recombine\_full"{]}}). Cycle 2 will wire
the remaining components :

\begin{itemize}
\tightlist
\item
  \textbf{α-stream raw traces} input channel (currently P\_max-only
  declared but not consumed) --- requires firehose ring buffer with
  bounded retention
\item
  \textbf{ATTENTION\_PRIOR canal-4} output channel --- requires the
  attention prior bounding invariant (S4) and downstream wiring to
  consumer modules
\item
  \textbf{\texttt{recombine\_full}} operation variant --- full VAE
  encoder / decoder pair beyond the C-Hobson light interpolation
  skeleton
\end{itemize}

With P\_max real-wired, hypothesis H2 (P\_max equivalence vs P\_equ
within ±5\%) becomes a real comparison rather than the cycle-1
self-equivalence smoke test.

\subsubsection{9.3 Real fMRI lab
partnership}\label{real-fmri-lab-partnership}

Cycle 1 locks Studyforrest as the fMRI fallback (G1 Branch A). Cycle 2
pursues active partnership with one or more fMRI labs identified via the
T-Col reviewer outreach :

\begin{itemize}
\tightlist
\item
  \textbf{Huth Lab} (UT Austin) : Narratives dataset
\item
  \textbf{Norman Lab} (Princeton) : episodic memory studies
\item
  \textbf{Gallant Lab} (UC Berkeley) : naturalistic stimulus-driven BOLD
\end{itemize}

A real lab partnership would enable RSA on \textbf{task-controlled}
linguistic stimuli rather than the narrative-comprehension fallback
Studyforrest provides. This would strengthen H3 (monotonic
representational alignment) which reached only borderline significance
in the cycle-1 synthetic pipeline validation (run\_id
\texttt{syn\_s15\_3\_g4\_synthetic\_pipeline\_v1}, dump
\texttt{docs/milestones/ablation-results.json}).

\subsubsection{9.4 Multi-substrate Conformance Criterion
validation}\label{multi-substrate-conformance-criterion-validation}

The strongest theoretical claim of Framework C-v0.5.0+STABLE ---
substrate-agnosticism via DR-3 Conformance Criterion --- needs empirical
validation across more than two substrates to be defensible at peer
review. Cycle 2 establishes the validation matrix : for each candidate
substrate (cycle-1 reference implementation ✅, E-SNN, hypothetical
transformer-based instance), verify all three conformance conditions
(signature typing, axiom property tests passing, BLOCKING invariants
enforceable).

A reusable conformance test suite (drafted in cycle 1
\texttt{tests/conformance/}) is the foundation. Cycle 2 will extend it
with substrate-specific adapters and run the full suite against each
candidate substrate, producing a conformance report publishable as a
supplementary artifact for Paper 1 (or as the main contribution of Paper
2's engineering ablation paper).

\begin{center}\rule{0.5\linewidth}{0.5pt}\end{center}

\subsection{10. References}\label{references}

→ See \texttt{references.bib} (16 entries cycle-1 stub, will extend to
\textasciitilde30-40 in S20-S22 as the full draft is rendered). BibTeX
integration via \texttt{\textbackslash{}bibliography\{references\}} in
the LaTeX render.

Key citations (alphabetical) : - Diekelmann \& Born 2010 (sleep memory)
- French 1999 (catastrophic forgetting) - Friston 2010 (FEP) - Hobson
2009 (REM dreaming) - Kirkpatrick 2017 (EWC) - McClelland 1995 (CLS) -
McCloskey \& Cohen 1989 (forgetting) - Rao \& Ballard 1999 (predictive
coding) - Rebuffi 2017 (iCaRL) - Shin 2017 (generative replay) - Solms
2021 (consciousness) - Stickgold 2005 (consolidation) - Tononi \&
Cirelli 2014 (SHY) - van de Ven 2020 (brain-inspired replay) - Walker \&
Stickgold 2004 (consolidation) - Whittington \& Bogacz 2017 (predictive
coding)

\begin{center}\rule{0.5\linewidth}{0.5pt}\end{center}

\subsection{Word count summary (target : \textasciitilde5000 words main
+ supp)}\label{word-count-summary-target-5000-words-main-supp}

{\def\LTcaptype{none} % do not increment counter
\begin{longtable}[]{@{}lll@{}}
\toprule\noalign{}
Section & Target & Status \\
\midrule\noalign{}
\endhead
\bottomrule\noalign{}
\endlastfoot
§1 Abstract & ≤250 & drafted (\textasciitilde265, needs trim) \\
§2 Introduction & ≤1500 & drafted (\textasciitilde1200) \\
§3 Background & ≤1500 & drafted (\textasciitilde1500) \\
§4 Framework & condensed in main + spec ref & done \\
§5 Implementation & condensed & done \\
§6 Methodology & ≤1500 & drafted (\textasciitilde1500) \\
§7 Results & ≤2000 & drafted (placeholder) \\
§8 Discussion & ≤1500 & drafted (\textasciitilde1500) \\
§9 Future Work & ≤700 & drafted (\textasciitilde700) \\
§10 References & n/a & 16 entries stub \\
\end{longtable}
}

\textbf{Estimated total} : \textasciitilde10000 words (needs aggressive
trim for Nature HB 5000-word main-text discipline ; supp can absorb
overflow).

\begin{center}\rule{0.5\linewidth}{0.5pt}\end{center}

\subsection{Notes for revision}\label{notes-for-revision}

\begin{itemize}
\tightlist
\item
  Render via Quarto / pandoc into PDF + LaTeX for arXiv submit (S21.1)
\item
  Insert OSF DOI in §6.1 once OSF lock is completed
\item
  Replace synthetic placeholders in §7 with real ablation values post
  S20+
\item
  Trim §1 abstract to ≤250 words
\item
  Trim §3 + §6 + §8 to fit overall main-text budget
\item
  Add Figures (1 architecture diagram, 2 results boxplot, 3 Jonckheere
  trend, 4 four-pillars conceptual)
\item
  Bibtex render with proper \texttt{\textbackslash{}cite\{\}} calls
\end{itemize}

\bibliography{../references}

\end{document}
